\documentclass[12pt,a4paper]{article}

\usepackage[utf8]{inputenc}
\usepackage[T1]{fontenc}
\usepackage{lmodern}
\usepackage[margin=1in]{geometry}
\usepackage{hyperref}
\usepackage{setspace}

\hypersetup{
    colorlinks=true,
    linkcolor=blue,
    citecolor=blue,
    urlcolor=blue
}

\title{\textbf{What the Supreme Court Tariff Ruling Really Means for Indonesia}}
\author{Teguh Wicaksono\thanks{The writer is a faculty member at the Faculty of Economics and Business, Universitas Islam Internasional Indonesia (UIII).}}
\date{}

\begin{document}

\maketitle
\thispagestyle{empty}

\onehalfspacing

The timing could not have been more dramatic. Since July 2025, when Trump threatened to reimpose 32 percent tariffs, Indonesia's negotiators worked under extraordinary pressure---months of shuttle diplomacy, concession after concession, culminating in the Agreement of Reciprocal Trade signed by Presidents Prabowo and Trump in Washington on February 19. Twenty-four hours later, the Supreme Court struck down the very legal authority---the International Emergency Economic Powers Act---under which that deal had been negotiated.

Indonesia's hard-won agreement was signed on a foundation that no longer exists.

\bigskip
\noindent\textbf{The numbers behind the disruption}

\medskip
Using 2024 trade data (\$28.05 billion in US imports from Indonesia) and import-weighted calculations, I estimate Indonesia's effective tariff rate across four scenarios.

Before the crisis, Indonesia faced an average effective tariff of 5.0 percent, driven by high MFN rates on textiles (12 percent) and footwear (11 percent). Liberation Day pushed that to 36.6 percent. The bilateral deal brought it to 19.7 percent, with zero-tariff exemptions for palm oil, rubber, coffee, cocoa, and semiconductors providing relief for \$6.1 billion in trade.

After the ruling, with the administration pivoting to a flat 10 percent tariff under Section 122 of the Trade Act of 1974, Indonesia's effective rate now stands at 15.3 percent---lower than the deal, but with a very different distribution of pain.

\bigskip
\noindent\textbf{A tale of two Indonesias}

\medskip
The shift creates winners and losers among Indonesian exporters, and the dividing line runs between manufacturing and commodities.

For Indonesia's textile, footwear, and furniture industries---which employ millions of workers and account for over \$6.3 billion in US-bound exports---the Supreme Court ruling is unexpectedly beneficial. Their tariffs drop from 29--30 percent under the deal to 21--22 percent under Section 122. Garment workers in Tangerang and shoe factory employees in Sidoarjo face marginally better prospects today than they did under the celebrated deal.

For commodity exporters, the picture inverts sharply. Palm oil producers, who enjoyed a hard-fought zero-tariff exemption under the deal, now face a flat 12 percent tariff (10 percent Section 122 plus 2 percent MFN). Rubber tappers, coffee growers, and cocoa farmers see similar jumps. The 1,819 exempted tariff lines---the crown jewel of Indonesia's negotiation---have evaporated overnight.

The net effect is a \$1.2 billion reduction in total duties compared to the deal scenario. But that aggregate number masks real losses for specific communities and value chains.

\bigskip
\noindent\textbf{What Indonesia conceded---and whether it still should}

\medskip
The more consequential question concerns the concessions Indonesia made to secure the deal. To get 19 percent and 1,819 zero-tariff lines, Jakarta---as this newspaper reported---``upended several policies it has maintained for years.'' The list is sobering: local content requirements (TKDN) exempted for US companies, halal certification relaxed for US goods, cross-border data transfer rights granted despite the US lacking federal data privacy law, critical mineral export restrictions loosened for American buyers, agricultural markets thrown open to 50,000 tons of US beef annually plus wheat, soybeans, and corn, and Freeport-McMoRan's Papua mining license extended 20 years to 2061. Add \$33 billion in purchase commitments and Article 5.1's requirement that Indonesia ``reflect US trade restrictions on third countries''---effectively a sanctions alignment clause.

These concessions were predicated on reciprocity. With access now governed by a blunt, non-negotiable 10 percent tariff rather than the carefully structured deal, the fundamental bargain has shifted. International trade law suggests that when one side's obligations are voided by judicial action, the other side's commitments warrant reassessment. Indonesia's parliamentary ratification---required within 90 days---should not proceed until this legal uncertainty resolves.

\bigskip
\noindent\textbf{The ticking clock}

\medskip
Section 122 carries a critical constraint: it expires after 150 days without congressional authorization, around July 20, 2026. If Congress does not act, the additional tariff disappears and Indonesia's effective rate reverts toward the 5 percent baseline.

This creates a peculiar landscape. The 10 percent applies to everyone equally---Indonesia, Vietnam, Bangladesh, the EU. Indonesia's 19 percent deal rate had actually disadvantaged it relative to competitors facing only the universal baseline; the Supreme Court has paradoxically leveled the playing field. But this uniformity also means the US has lost its primary tool for bilateral negotiations.

\bigskip
\noindent\textbf{What Indonesia should do now}

\medskip
First, pause. Do not rush to implement concessions until their legal basis is clarified. Prudence, not provocation, should guide the next steps.

Second, pursue refund claims. The 19 percent reciprocal tariff paid from August 2025 through February 2026 was collected under now-invalid authority. These duties---potentially over \$500 million---should be recoverable through US Customs protests.

Third, use the 150-day window. With all competitors facing the same rate, Indonesian manufacturers should accelerate shipments to regain lost market share.

Fourth, diversify. The past year has demonstrated that dependence on any single market creates existential vulnerability. RCEP, the EU--Indonesia CEPA, and deeper engagement with Middle Eastern and African markets deserve renewed urgency.

The Supreme Court did not deliver Indonesia a gift. It delivered complexity---a lower average tariff but lost exemptions, a level playing field but a temporary one, reduced US leverage but also reduced Indonesian bargaining power. Navigating this demands clear-eyed analysis and strategic patience.

The bilateral relationship between the world's third- and fourth-largest democracies is too important for either side to let it be defined by the next 150 days.

\end{document}
